Este capítulo explica los conceptos, teorías y definiciones necesarias para entender cómo se construyó Pymetory. Está organizado en tres bloques.

El \textbf{Marco Teórico} expone los fundamentos que sustentan el sistema: desde los principios de gestión de inventarios y las estrategias de rotación de productos (FEFO, FIFO, LIFO) hasta la arquitectura RAG que hace posible el módulo de consulta inteligente, pasando por las tecnologías de desarrollo web que se usaron (Laravel con patrón MVC, React con Inertia.js).

El \textbf{Estado del Arte} analiza las soluciones comerciales y académicas que ya existen, identificando qué hacen bien y qué brechas deja sin cubrir este proyecto. Esto permite entender por qué Pymetory no es ``otro sistema de inventarios más'' sino una propuesta que cierra vacíos concretos.

El \textbf{Marco Conceptual} presenta un glosario de términos técnicos y de negocio, todos contextualizados en el caso real de la panificadora en Santiago de Cali que sirvió como laboratorio de requisitos. Si en algún momento aparece un acrónimo o concepto que no queda claro, este glosario es el lugar para buscarlo.
